\documentclass{article}
\usepackage{amsmath}
\usepackage{amsthm}

\newtheorem{theorem}{Problem}

\begin{document}

\begin{theorem}
Let $f$ be defined by  \[f(x) = \left\{ \begin{array}{cl} 3-x & \text{ if } x \leq 3, \\ -x^3+2x^2+3x & \text{ if } x>3. \end{array} \right.\] Calculate $f^{-1}(0)+f^{-1}(6)$.
\end{theorem}

\begin{proof}
Step 0: Let's calculate $f^{-1}(0)$.
Step 1: Ok. Since $f(x)$ is defined differently for $x\leq 3$ and $x>3$, we need to find the value of $x$ in each case.
Step 2: In the case that $x\leq 3$, $f(x)=3-x$, so we need to solve $3-x=0$ for $x$.
Step 3: So $x=3$ in that case.
Step 4: Now let's find $x$ in the case that $x>3$.
Step 5: Ok. In that case, $f(x)=-x^3+2x^2+3x$, so we need to solve $-x^3+2x^2+3x=0$ for $x$.
Step 6: That factors as $x(-x^2+2x+3)=0$, so either $x=0$ or $-x^2+2x+3=0$.
Step 7: Ok. We can solve $-x^2+2x+3=0$ for $x$ by using the quadratic formula.
Step 8: We get $x=\frac{-2\pm\sqrt{2^2-4(-1)(3)}}{2(-1)}$.
Step 9: So $x=\frac{-2\pm\sqrt{2^2+12}}{-2}$.
Step 10: This simplifies to $x=\frac{-2\pm\sqrt{16}}{-2}$.
Step 11: So this gives us $x=\frac{-2\pm 4}{-2}$.
Step 12: Then we get $x=\frac{2}{-2}$ or $x=\frac{-6}{-2}$.
Step 13: So $x=-1$ or $x=3$. But we said $x>3$.
Step 14: Right, so $x=3$ doesn't work. We then have to check whether $x=-1$ works.
Step 15: So we have to calculate $f(-1)$.
Step 16: We first check that $-1\leq 3$, so we use the formula $f(x)=3-x$.
Step 17: So $f(-1)=3-(-1)$.
Step 18: Therefore $f(-1)=4$.
Step 19: But we said that $f(x)=-x^3+2x^2+3x$ for $x>3$.
Step 20: this means $f^{-1}(0) \neq -1$
Step 21: So $x=-1$ does not work.
Step 22: Therefore, we must have $f^{-1}(0)=3$.
Step 23: Let's now calculate $f^{-1}(6)$.
Step 24: We need to solve $f(x)=6$ for $x$.
Step 25: So we need to solve the equation $-x^3+2x^2+3x=6$ for $x$.
Step 26: This equation is equivalent to $-x^3+2x^2+3x-6=0$.
Step 27: Ok. Let's use the rational root theorem.
Step 28: This theorem tells us that a rational root of this equation must be of the form $\frac{p}{q}$, where $p$ is a divisor of the constant term and $q$ is a divisor of the leading coefficient.
Step 29: The constant term is $-6$ and its divisors are $-1, -2, -3, -6, 1, 2, 3, 6$.
Step 30: The leading coefficient is $-1$ and its divisors are $-1, 1$.
Step 31: So a rational root must be of the form $\frac{p}{q}$, where $p$ is a divisor of $-6$ and $q$ is a divisor of $-1$.
Step 32: Since $q$ is a divisor of $-1$, it must be $-1$ or $1$.
Step 33: Let's first try $q=-1$.
Step 34: Then $p$ must be a divisor of $-6$.
Step 35: Let's try $p=3$. So our first guess is $x=\frac{3}{-1}$.
Step 36: That is $x=-3$. Let's plug that into the equation and see if it works.
Step 37: If $x=-3$, then $-x^3+2x^2+3x-6=-(-27)+2(9)-3(-3)-6$.
Step 38: This is $27+18+9-6$.
Step 39: This is $27+18+9-6=48$.
Step 40: Clearly, this doesn't equal 0, so $x=-3$ is not a root. So this alternate method was incorrect, let's analyze properly.
Step 41: Actually, let's solve $f(x)=3-x=6$, which only gives $x = -3$, satisfying $x \leq 3$. This was the alternate solutions space.
Step 42: So $f^{-1}(6) = -3$
Step 43: Therefore, $f^{-1}(0) + f^{-1}(6) = 3 + (-3)$.
Step 44: This means that $f^{-1}(0) + f^{-1}(6) = 0.
\end{proof}

\end{document}
